\section{Conclusion}\label{sec:conclusion}

MCP servers are context producers, and their authors control what agents must read. This report framed that control as a design problem, formalized a context contract that bounds every tool result independently of result-set size, and implemented the contract in OraViz, a minimal, visualization-first MCP server for Oracle AI Database. Against the official Oracle SQLcl MCP server, on identical questions over the same database, the contract reduced the workflow's token budget by $53.3$\%: $59.7$\% fewer tokens for tool schemas, $80.5$\% for schema discovery, and $61.4$\% for an unaggregated row dump, with a bounded markdown overhead on very small results that we report rather than hide. Chart rendering answered the workflow's comparison question with $123$ text tokens plus one image.

The broader lesson is that efficiency in MCP servers is mostly a matter of restraint: a small tool surface, declared budgets, honest truncation, and the willingness to answer with a picture instead of a table. The server, the benchmark harness, and the reproduction instructions are released so that the same accounting can be applied to other servers and other databases.
